\uuid{VIj3}
\titre{Ensemble}
\auteur{bodin ; hanani ; mzari}
\organisation{exo7}
\licence{CC BY-NC-SA 4.0}
\datecreate{2020-01-20}
\niveau{L1}
\module{Algèbre}
\chapitre{Logique, ensemble, raisonnement}
\sousChapitre{Ensemble}
\difficulte{2}
\theme{Ensembles, applications, ensemble}
\exo7qcmid{lille-1.ensembles.012}

\contenu{
\begin{question}
\qclassification{100.02,101.01,102.02,102.02}
\qtags{motcle=ensemble}
Soit $A$ une partie d'un ensemble $E$ telle que $A\neq E$. On note $\overline{A}$ le complémentaire de $A$ dans $E$. Quelles sont les bonnes réponses ?
\begin{answers}
    \bad{$A\cap \overline{A}=E$}
    \good{$A\cap \overline{A}=\varnothing$}
    \good{$A\cup\overline{A}=E$}
    \bad{$A\cup \overline{A}=A$}
\end{answers}
\begin{explanations}
S'il existe $x\in E$ tel que $x\in A\cap \overline{A}$ alors $(x\in A$ et $x\notin A)$. Ceci est absurde. Donc $A\cap \overline{A}=\varnothing$. De même $x\in E\Rightarrow (x\in A$ ou $x\notin A)$. Donc que $E\subset A\cup\overline{A}\subset E$.
\end{explanations}
\end{question}
}
